% SGA 3, Exposé IX, tradução brasileira.
% Escopo: §4 desde o início até a primeira metade do Lema 4.4.
% Autoridade: Polo--Gille Exp9-8nov09.pdf, pp. locais 10--14.

\setcounter{footnote}{10}

\SGARefTarget{sga3:ix:section:4}\section*{4. O teorema de densidade}
\addcontentsline{toc}{section}{4. O teorema de densidade}

Este teorema (cf. \SGARefLink{sga3:ix:theorem:4:7}{4.7 abaixo})%
\begingroup
\renewcommand{\thefootnote}{\(\ast\)}%
\footnote{Todos os resultados desde
\SGARefLink{sga3:ix:definition:4:1}{4.1} até
\SGARefLink{sga3:ix:corollary:4:6}{4.6} estão contidos em EGA
IV\(_3\), 11.10, ao qual remetemos o leitor para uma exposição sistemática
da noção de densidade esquemática.}%
\addtocounter{footnote}{-1}%
\endgroup,
juntamente com o teorema\footnote{Nota do editor: o teorema de
«algebricidade» \SGARefLink{sga3:ix:theorem:7:1}{7.1}.} da
\SGARefLink{sga3:ix:section:7}{Seção~7}, será a ferramenta essencial na
presente Exposição e nas duas seguintes para passar das
propriedades infinitesimais dos grupos de tipo multiplicativo
desenvolvidas acima às suas propriedades «finitas».

\medskip
\noindent\SGARefTarget{sga3:ix:definition:4:1}\textbf{Definição 4.1.}
\textit{Seja \(X\) um pré-esquema. Diz-se que uma família
\((Z_i)_{i\in I}\) de subpré-esquemas de \(X\) é \emph{esquematicamente
densa} se, para todo aberto \(U\) de \(X\) e todo subpré-esquema fechado
\(U_0\) de \(U\) que contenha todos os \(Z_i\cap U\), tem-se
\(U_0=U\).}

\textit{Diz-se que um subpré-esquema \(Z\) de \(X\) é
\emph{esquematicamente denso em \(X\)} se a família de um só elemento
formada por \(Z\) for esquematicamente densa.}

É imediato (cf. EGA IV\(_3\), 11.10.1) que a definição equivale a
dizer que, para todo aberto \(U\) de \(X\), toda seção \(f\) de
\(\mathscr O_U\) que se anule sobre cada \(Z_i\cap U\) é nula. De maneira
equivalente, a interseção dos núcleos dos homomorfismos canônicos
\SGARefTarget{sga3:ix:definition:4:1:display:01}
\[
  u_i:\mathscr O_X\longrightarrow (v_i)_*(\mathscr O_{Z_i})
\]
é nula, onde \(v_i:Z_i\to X\) é a imersão canônica.%
\footnote{Nota do editor: todos os resultados sobre densidade esquemática
enunciados em EGA IV\(_3\), \S11.10 deduzem-se dos resultados sobre
«famílias separadoras de homomorfismos» provados em loc.\ cit., \S11.9,
ao qual remeterão algumas notas editoriais posteriores.}
Quando \(X\) está sobre um pré-esquema \(S\), isso equivale também à
afirmação seguinte: para todo aberto \(U\) de \(X\), todo
\(S\)-pré-esquema \(Y\) separado sobre \(S\) e todo par de morfismos
\(u,v:U\to Y\) que coincidam sobre os \(Z_i\cap U\), tem-se \(u=v\).
Com efeito, \(u=v\) equivale a \(U_0=U\), onde \(U_0\) é a imagem
inversa da diagonal em \(Y\times_S Y\) pelo \(S\)-morfismo
\(U\to Y\times_S Y\) definido por \((u,v)\). A diagonal é um
subpré-esquema fechado de \(Y\times_S Y\), de modo que \(U_0\) é um
subpré-esquema fechado de \(U\) que contém todos os \(Z_i\cap U\). A
densidade esquemática dá, pois, \(U_0=U\), e portanto \(u=v\). Para a
recíproca, tome-se \(Y=\operatorname{Spec}\mathscr O_S[T]\).

Na terminologia introduzida ao final de EGA I, 9.5, dizer que o
subpré-esquema \(Z\) é esquematicamente denso em \(X\) significa também
que \(X\) é o \emph{subpré-esquema aderência} de \(Z\) em \(X\).

% Página-fonte: Exp9 p. local 11; leitor completo p. 657.
\medskip
\noindent\SGARefTarget{sga3:ix:lemma:4:2}\textbf{Lema 4.2.}
\textit{Seja \(X\) um \(S\)-pré-esquema plano e seja \((Z_i)_{i\in I}\) uma
família de subpré-esquemas de \(X\), planos sobre \(S\). Seja \(S_0\) um
subpré-esquema de \(S\) definido por um ideal nilpotente \(\mathscr J\), e
suponhamos que os módulos \(\mathscr J^n/\mathscr J^{n+1}\) sejam localmente
livres sobre \(S_0\). Sejam \(X_0,Z_{i0}\) os obtidos de \(X,Z_i\) pela
mudança de base \(S_0\to S\). Se \((Z_{i0})_{i\in I}\) for
esquematicamente densa em \(X_0\), então \((Z_i)_{i\in I}\) será
esquematicamente densa em \(X\).}

\emph{Demonstração.}
Suponhamos \(\mathscr J^{n+1}=0\), onde \(n\geq0\), e raciocinemos por
indução sobre \(n\); a afirmação é trivial para \(n=0\). Definamos
\(S_m,X_m,Z_{im}\) por redução módulo \(\mathscr J^{m+1}\), como de
costume. A hipótese de indução já implica que
\((Z_{i,n-1})_{i\in I}\) é esquematicamente densa em \(X_{n-1}\).
Após substituir \(X\) por um aberto induzido, basta provar que toda
seção \(f\) de \(\mathscr O_X\) que se anule sobre todos os \(Z_i\) é
nula.

A seção \(f_{n-1}\) de
\[
  \mathscr O_{X_{n-1}}
  =\mathscr O_X/\mathscr J^n\mathscr O_X
\]
definida por \(f\) anula-se sobre os \(Z_{i,n-1}\), logo é nula. Por
conseguinte, \(f\) é uma seção de \(\mathscr J^n\mathscr O_X\). Como
\(X\) é plano sobre \(S\), existe um isomorfismo
\SGARefTarget{sga3:ix:lemma:4:2:display:01}
\[
  \mathscr J^n\mathscr O_X
  \xleftarrow{\ \sim\ }
  \mathscr E\otimes_{\mathscr O_{S_0}}\mathscr O_{X_0},
  \qquad
  \mathscr E=\mathscr J^n=\mathscr J^n/\mathscr J^{n+1}.
\]
Do mesmo modo, como cada \(Z_i\) é plano sobre \(S\), a
restrição \(f_i\) de \(f\) a \(Z_i\) pode ser vista como uma seção de
\SGARefTarget{sga3:ix:lemma:4:2:display:02}
\[
  \mathscr J^n\mathscr O_{Z_i}
  \xleftarrow{\ \sim\ }
  \mathscr E\otimes_{\mathscr O_{S_0}}\mathscr O_{Z_{i0}}.
\]
Por hipótese, todos os \(f_i\) são nulos. O módulo \(\mathscr E\) é
localmente livre por hipótese e, portanto, também o é
\[
  \mathscr F
  =\mathscr E\otimes_{\mathscr O_{S_0}}\mathscr O_{X_0}.
\]
Assim, \(f\) é uma seção do módulo localmente livre \(\mathscr F\) sobre
\(X_0\) cuja restrição a cada \(Z_{i0}\) é nula. Como a família
\((Z_{i0})_{i\in I}\) é esquematicamente densa em \(X_0\), segue-se
imediatamente que \(f=0\). \qed

% Página-fonte: Exp9 p. local 12; leitor completo p. 658.
\medskip
\noindent\SGARefTarget{sga3:ix:corollary:4:3}\textbf{Corolário 4.3.}
\textit{Seja \(X\) um \(S\)-pré-esquema localmente noetheriano, plano sobre
\(S\), e seja \((Z_i)_{i\in I}\) uma família de subpré-esquemas de \(X\),
planos sobre \(S\). Suponhamos que, para todo \(s\in S\), a família de
fibras \((Z_{is})_{i\in I}\) seja esquematicamente densa em \(X_s\).
Então \((Z_i)_{i\in I}\) é esquematicamente densa em \(X\).}

\emph{Demonstração.}
Após substituir \(X\) por um aberto induzido, basta provar que toda
seção \(f\) de \(\mathscr O_X\) cuja restrição a todos os \(Z_i\) é
nula é ela própria nula.%
\footnote{Nota do editor: desenvolveu-se o original no que segue.}
Para isso, basta mostrar que a imagem de \(f\) em \(\mathscr O_{X,x}\)
é nula para todo \(x\in X\). Seja \(\mathfrak m_x\) o ideal maximal de
\(\mathscr O_{X,x}\), e seja \(\mathfrak m_s\) o de
\(\mathscr O_{S,s}\), onde \(s\) é a imagem de \(x\) em \(S\). Como
\(\mathscr O_{X,x}\) é noetheriano,
\[
  \bigcap_{n\in\mathbb N}\mathfrak m_x^n=0,
  \qquad\text{y, con mayor razón,}\qquad
  \bigcap_{n\in\mathbb N}\mathscr O_{X,x}\mathfrak m_s^n=0.
\]
Basta, pois, mostrar, para todo inteiro \(n\), que a seção induzida por
\(f\) sobre
\[
  X\times_S
  \operatorname{Spec}\bigl(\mathscr O_{S,s}/\mathfrak m_s^{n+1}\bigr)
\]
é nula. Por hipótese,
\[
  Z_{is}=Z_i\otimes_{\mathscr O_{S,s}}\kappa(s)
\]
é uma família esquematicamente densa em
\[
  X_s=X\otimes_{\mathscr O_{S,s}}\kappa(s).
\]
Reduzimo-nos, assim, ao caso em que \(S\) é o espectro de um anel local
cujo ideal maximal \(\mathscr J\) é nilpotente,
\(S_0=\operatorname{Spec}\kappa(s)\), e se cumprem as hipóteses do
\SGARefLink{sga3:ix:lemma:4:2}{Lema~4.2}. Segue-se a conclusão. \qed

\medskip
\noindent\SGARefTarget{sga3:ix:lemma:4:4}\textbf{Lema 4.4.}
\textit{Seja \(S\) um pré-esquema localmente noetheriano, seja \(X\) um
\(S\)-pré-esquema localmente noetheriano, plano sobre \(S\), e seja
\((Z_i)_{i\in I}\) uma família de subpré-esquemas de \(X\), planos sobre
\(S\). Suponhamos que, para todo \(s\in S\), a família de fibras
\((Z_{is})_{i\in I}\) seja esquematicamente densa em \(X_s\). Então,
para toda mudança de base \(S'\to S\), sem exigir que \(S'\) seja localmente
noetheriano, a família \((Z_i')_{i\in I}\) é esquematicamente densa em
\(X'\). Dito de outro modo, \((Z_i)_{i\in I}\) é universalmente
esquematicamente densa em \(X\) com respeito a \(S\) (cf. EGA IV\(_3\),
Definição 11.10.8).}

\emph{Demonstração.}%
\footnote{Nota do editor: para outra demonstração, que reduz ao caso em
que \(S'\) é localmente noetheriano e utiliza
\SGARefLink{sga3:ix:corollary:4:3}{4.3} e
\SGARefLink{sga3:ix:lemma:4:5}{4.5} como aqui, vejam-se EGA IV\(_3\),
11.9.16 e 11.9.12. Na última linha da demonstração de 11.9.16,
substitua-se 11.9.5 por 11.9.12.}
Seja \(f'\) uma seção de \(\mathscr O_{X'}\) sobre um aberto \(U'\) de
\(X'\), que se anule sobre os \(Z_i'\); devemos provar que \(f'=0\).
Fixemos \(s'\in S'\). Basta provar que \(f'\) se anula em todo ponto de
\(U'\) situado sobre \(s'\). Se \(s\) é a imagem de \(s'\) em \(S\),
após substituir \(S,S'\) pelos espectros de seus anéis locais em
\(s,s'\), podemos supor que \(S,S'\) são locais e que \(s,s'\) são seus
pontos fechados. Podemos supor também que \(X\) é afim, de modo que
\SGARefTarget{sga3:ix:lemma:4:4:display:01}
\[
  \begin{gathered}
    S=\operatorname{Spec}(A),\qquad
    S'=\operatorname{Spec}(A'),\\
    X=\operatorname{Spec}(B),\qquad
    X'=\operatorname{Spec}(B'),\qquad
    B'=B\otimes_A A',
  \end{gathered}
\]
onde \(A,A'\) são locais, \(A\to A'\) é um homomorfismo local, e
\(A,B\) são noetherianos. Podemos supor ainda que \(U'\) é afim, da
forma \(X'_{g'}\), com \(g'\in B'\).

Para toda \(A\)-subálgebra \(A''\) de \(A'\), ponhamos
\[
  S''=\operatorname{Spec}(A'')
\]
e sejam \(X'',Z_i''\) os obtidos de \(X,Z_i\) pela mudança de base
\(S''\to S\). Obtém-se o diagrama
\SGARefTarget{sga3:ix:lemma:4:4:diagram:01}
\[
\begin{tikzcd}[column sep=large,row sep=large]
  % La fuente imprime visiblemente subíndices n en los nodos superiores
  % central y derecho; se conservan pese a la prosa circundante.
  Z^i \arrow[d] & Z_n^{i''} \arrow[l]\arrow[d]
    & Z_n^{i'} \arrow[l]\arrow[d] \\
  X \arrow[d] & X'' \arrow[l]\arrow[d] & X' \arrow[l]\arrow[d] \\
  S & S'' \arrow[l] & S' \arrow[l].
\end{tikzcd}
\]

% Página-fonte: Exp9 p. local 13; leitor completo p. 659.
No que segue, limitamo-nos aos \(A''\) que são localizações de
\(A\)-subálgebras \(A_1''\) de \(A'\), de tipo finito sobre \(A\), em
\(\mathfrak m'\cap A_1''\), onde \(\mathfrak m'\) é o ideal maximal de
\(A'\). Assim, os homomorfismos \(A\to A''\to A'\) são locais. Esses
\(A''\) formam uma família filtrante crescente de subálgebras cuja união,
e portanto cujo limite direto, é \(A'\). Do mesmo modo,
\[
  B'=\varinjlim_{A''}\bigl(B\otimes_A A''\bigr).
\]
Existem, pois, tal \(A''\) e um elemento
\(g''\in B''=B\otimes_A A''\) cuja imagem em \(B'\) é \(g'\).

Pelo \SGARefLink{sga3:ix:lemma:4:5}{Lema~4.5 abaixo}, a propriedade de
\((Z_{is})_{i\in I}\) ser esquematicamente densa em \(X_s\)
conserva-se por toda mudança de base \(S''\to S\).%
\footnote{Nota do editor: acrescentou-se a frase precedente e desenvolveu-se
o original no que segue.}
Por conseguinte, como cada \(S''\) é localmente noetheriano, substituir
\(S\) por um \(S''\) adequado e \(X,Z_i\) por \(X'',Z_i''\) conserva as
hipóteses do \SGARefLink{sga3:ix:lemma:4:4}{Lema~4.4}.

Podemos supor, portanto, que \(g'\) provém de um elemento \(g\in B\).
Após substituir \(X\) pelo aberto \(X_g=U\), podemos supor que
\(U'=X'\). Do mesmo modo, podemos supor que \(f'\) provém de uma
seção \(f\) de \(\mathscr O_X\) sobre \(X\).

Seja \(Y=V(f)\) o subesquema de \(X\) definido por \(f\), e ponhamos
\[
  Y_i=Z_i\times_XY.
\]
Assim, \(Y_i\) é o subesquema fechado \(V(f_i)\) de \(Z_i\), onde
\(f_i\) é a seção de \(\mathscr O_{Z_i}\) induzida por \(f\).
Denotando por \(Y',Y_i'\) os \(S'\)-esquemas obtidos por mudança de base,
tem-se
\SGARefTarget{sga3:ix:lemma:4:4:display:02}
\[
  Y'=V(f'),
  \qquad
  Y_i'=Z_i'\times_{X'}Y'=V(f_i'),
\]
e valem relações análogas para \(Y'',Y_i''\). A afirmação de que
\(f'\) se anula sobre os \(Z_i'\) equivale a
\[
  Y_i'=Z_i'\qquad\text{para todo }i.
\]
\footnote{Nota do editor: o original acrescentava «ao menos em todo ponto
\(x'\) de \(X'\) situado sobre \(s'\)», mas essa restrição parece
desnecessária.}

Seja \(\mathfrak m\) o ideal maximal de \(A\). Para todo \(n\geq0\),
introduzamos o subesquema fechado
\[
  S_n=\operatorname{Spec}(A/\mathfrak m^{n+1})
\]
de \(S\), e sejam \(X_n,Y_n,Z_n^i,Y_n^i\) os obtidos de
\(X,Y,Z_i,Y_i\) pela mudança de base \(S_n\to S\). Mais geralmente,
para todo \(S\)-pré-esquema \(P\), escrevamos \(P_n=P\times_SS_n\).

Para todo \(n\) e \(i\in I\), consideremos o funtor
\SGARefTarget{sga3:ix:lemma:4:4:display:03}
\[
  F_n^i
  =\prod_{Z_n^i/S_n}Y_n^i/Z_n^i:
  (\mathbf{Sch})_{/S_n}^{\circ}\longrightarrow\mathbf{Set},
\]
definido, para todo \(P\) sobre \(S_n\), por
\SGARefTarget{sga3:ix:lemma:4:4:display:04}
\[
F_n^i(P)
=\Gamma\bigl((Y_n^i)_P/(Z_n^i)_P\bigr)
=
\begin{cases}
  \varnothing,&(Y_n^i)_P\ne(Z_n^i)_P,\\
  \{\operatorname{id}\},&(Y_n^i)_P=(Z_n^i)_P.
\end{cases}
\]
Como \(S_n\) é local artiniano e \(Z_n^i\) é plano, e portanto
essencialmente livre, sobre \(S_n\),
\SGARefLink{sga3:viii:theorem:6:4}{VIII, 6.4} mostra que cada \(F_n^i\)
é representado por um subpré-esquema fechado \(T_n^i\) de \(S_n\).

% Página-fonte: Exp9 p. local 14; leitor completo p. 660.
O completamento \(\widehat A\) do anel local \(A\) é noetheriano
porque \(A\) o é, e é o limite projetivo dos anéis \(A_n\).%
\footnote{Nota do editor: desenvolveu-se e corrigiu-se o original para
mostrar que o limite projetivo dos anéis \(\mathscr O(T_n^i)\)
define um subesquema formal fechado de \(\operatorname{Spf}(\widehat A)\).}
Seja \(K_n^i\) o ideal de
\[
  A_n=A/\mathfrak m^{n+1}
\]
que define \(T_n^i\), e seja \(K^i=\varprojlim_nK_n^i\); este é um ideal
de \(\widehat A\).
